\section{Introduction}
\label{sec:introduction}

% Author text.

Fertility is falling almost everywhere. In the United States, the period
fertility rate fell from about 2.1 children per woman in 2007 to 1.62 in
2023, and the world as a whole is likely below replacement for the first
time.\footnote{The total fertility rate was 2.12 in 2007
\citep{cboDemographicOutlook2026} and 1.621 in 2023
\citep{ostermanBirthsFinal2023}. On the world as a whole, see
\citet{fernandezVillaverdeNorrickTerraIncognita2026}.} The Congressional
Budget Office now projects that deaths will exceed births from 2030, and
governments have started to respond, often through housing. Hungary writes
off part of a family's home loan at each birth from the second child, and
South Korea lends at subsidized rates to households with a
newborn.\footnote{Hungary writes off 10 million forints of a 3\% family home
loan at the birth of each child from the second \citep{hungaryCSOKPlus}, and
South Korea lends up to 500 million won to households with a newborn
\citep{koreaNewbornLoan}. Poland repays 20,000 zloty of a family home loan at
a second birth and 60,000 zloty at a third \citep{polandFamilyRepayment}, and
Singapore sets aside up to 40\% of new public flats for couples with children
\citep{singaporeFPPS}. In the United States, the 2025 budget law raised the
child tax credit to \$2,200 and created accounts with a \$1,000 federal
deposit for children born from 2025 to 2028 \citep{publicLaw119-21}. In a
fourteen-country survey, housing is the third most cited reason for having
fewer children than desired \citep{unfpaSWP2025}.} In the United States, both
2024 presidential campaigns pointed to the cost of housing for young
families, and one proposed up to \$25,000 in down-payment
assistance.\footnote{See \citet{harrisCampaign2024}. The Republican
vice-presidential candidate said that young people are ``delaying having
families'' because ``they can't afford a place to raise a family''
\citep{vanceNewsNation2024}.} The logic is intuitive. Children need space,
and young households, who have little savings, find that space costly to pay
for.

Yet it is far from obvious which obstacle housing policy should remove. A
young couple that wants a child faces two costs. The first is the extra
space a family occupies every year its children live at home, whether it
rents or owns. The second, for a couple that buys, is the down payment needed
to obtain that space. Large homes are mostly owner-occupied: in 2023, 43\% of
owned homes in the United States had seven or more rooms, against 7.5\% of
rented homes.\footnote{Author's calculations from the 2023 American Housing
Survey.} If the down payment is what stands between young couples and a
family-size home, easier credit could raise births substantially. If instead
the recurring cost of space is what holds them back, easier credit may mainly
change who owns a home and when, and leave family size unchanged.

Two further features make the answer hard to read off the data. A child
commits a household to paying for its space for many years, including years
in which earnings fall, so policies that make this commitment safer may
matter as much as policies that make it cheaper. And births today become
households a generation later. As births decline, fewer families compete for
homes, and a policy that raises fertility raises the future demand for
housing and the price of the space the next generation of parents needs. The
effect of a policy on impact can therefore be a poor guide to its effect over
time. This paper asks which housing policies raise births, by how much, and
over what horizon.


This paper contributes to four literatures: on housing costs and fertility, on credit and fertility, on earnings risk and fertility, and on housing over the life cycle.

\paragraph{Housing costs and fertility.}
A large empirical literature finds that housing conditions shape childbearing. Fertility is lower where rent per room is higher \citep{simonTamuraHigherRents2009}, dwelling size predicts first births \citep{stromHousingFirstBirths2010}, and higher house prices raise births among owners but lower them among renters \citep{dettlingHousePricesBirth2014,lovenheimFamilyWealthShocks2013}. Other work finds that the cost of children has not risen: the price of a bundle of goods and services for children has not outpaced adult prices, and children require little more than an extra bedroom \citep{pattersonSarsons2026}, and state-level changes in rents do not line up with changes in birth rates \citep{kearneyPuzzleFallingUS2022}. I take the decline in births as given, and ask how the cost of family space shapes the response of fertility. Two recent papers model housing and fertility jointly. \citet{khorunzhinaMillerAmericanDream2022} estimate a dynamic model of female labor supply, fertility and the first home purchase at given house prices, and \citet{couillardBuildBabyBuild2025} estimates a model of fertility, living arrangements and location in which rents are the cost of housing. Relative to these papers, I study fertility together with saving, tenure and the choice of space under uninsurable earnings risk, in a housing market where the price is determined in equilibrium. This lets me separate the recurring cost of family space from the equity needed to buy it. I find that births respond to the first, not the second.

\paragraph{Credit and fertility.}
Quasi-experimental studies find that credit raises births. Households exposed to the deregulation of home lending are more likely to buy a home and have a child \citep{hacamoBabiesMortgageDeregulation2021}, lower interest rates on home loans raise births among households whose rate resets \citep{cummingMonetaryPolicyBirth2024}, federally backed home loans raised births and completed fertility during the baby boom \citep{dettlingKearney2026}, and winners of housing-credit lotteries in Brazil have more children, especially at ages 20 to 25 \citep{vandoornikFazioRamadoraiSkrastinsHousingFertility2025}. In these studies credit comes together with a home. A quantitative literature finds that credit relaxation moves ownership but not the size of the homes households live in or house prices \citep{kaplanMitmanViolanteHousing2020}, and that removing the down payment requirement does little to make ownership attractive \citep{halketVasudevSaving2014}. I reconcile the two. In my model, relaxing credit for everyone lets young households buy the small homes they would otherwise rent, which brings births forward and slightly lowers the completed fertility of the cohort now entering adulthood. Credit tied to children lowers the cost of the space a child requires, and raises completed fertility. Credit therefore matters for fertility only to the extent that it lowers the cost of family space.

\paragraph{Earnings risk and fertility.}
Because a child is an irreversible commitment, prudent parents who face more risk want fewer children \citep{fraserIncomeRisk2001}. Quantitative life-cycle models find that uninsurable earnings risk delays and reduces births \citep{sommerFertilityChoiceLife2016,ejrnaesJorgensenFamilyPlanning2020,nakakuniFertilityIncomeUncertainty2026}, and \citet{deLaCroixPommeretChildbearing2021} show that risk gives the postponement of a first birth an option value. Quasi-experimental studies find that job displacement and job insecurity lower fertility \citep{delBonoWeberWinterEbmer2012,clarkLepinteurJobInsecurity2022}. None of these papers includes housing. In my model, part of the commitment a child entails is the space it needs, paid for every period at a price set in the housing market. The response of fertility to earnings risk therefore depends on housing costs and on the down payment required to own family space.

\paragraph{Housing over the life cycle.}
Quantitative models of housing study tenure choice under earnings risk, down payments and transaction costs \citep{fisherGervaisHomeOwnership2011,halketVasudevSaving2014,boarGoreaMidriganLiquidity2022}, and the joint determination of house prices, rents and ownership \citep{kaplanMitmanViolanteHousing2020,greenwaldGurenCredit2025}. In this literature, rising earnings risk accounts for much of the decline in homeownership among young households \citep{fisherGervaisHomeOwnership2011,pazPardoHomeownership2024}. Household size, however, is exogenous or absent. I make fertility a choice. A birth raises the household's demand for space over its life cycle and, a generation later, the number of households that compete for housing. Housing policy therefore shapes both family size and the future demand for housing.
